\chapter{Justificativa da Escolha Metodológica}
\label{cap:justificativa-metodologica}

\section{Por que recombinação de fragmentos e árvores}

A recombinação de fragmentos sobre um catálogo finito garante, por construção, que cada molécula possa ser rastreada até os anéis, as pontes e os substituintes de origem, e dispensa bases de treinamento. Em contrapartida, a cobertura do espaço químico fica limitada ao catálogo, uma limitação reconhecida frente a modelos generativos treinados, que podem, em princípio, propor estruturas fora dele.

\section{Por que a TED e não outras medidas de distância}

A medida de distância define quais pares de anéis são considerados distantes e, portanto, quais combinações são priorizadas. As alternativas mais evidentes são as seguintes.

\begin{itemize}
	\item \textbf{Similaridade de Tanimoto sobre \textit{fingerprints} (ECFP).} É barata e muito usada, mas é calculada sobre a presença de subestruturas locais e não sobre a organização entre elas. Neste trabalho, a representação por árvore já separa explicitamente esqueleto, posição e fragmento, e a TED registra diferenças em cada um desses níveis. O Tanimoto sobre ECFP continua sendo a métrica adequada para avaliar a diversidade da molécula final (Capítulo~\ref{cap:objetivos}) e, por ser independente da árvore, evita avaliar o método com o mesmo critério que o orienta.
	\item \textbf{Distância de edição entre grafos (GED).} Opera diretamente sobre o grafo molecular, mas é NP-difícil em geral, o que a torna pouco prática para comparar muitos pares de anéis a cada geração.
	\item \textbf{TED para árvores ordenadas (Zhang--Shasha).} Tem algoritmo polinomial \cite{zhang_shasha_1989,bille_2005}, ao contrário da versão para árvores não ordenadas, que é NP-difícil \cite{zhang_statman_shasha_1992}. O custo é a necessidade de uma ordenação canônica dos filhos de cada nó, tratada no Capítulo~\ref{cap:metodologia}.
\end{itemize}

A escolha do algoritmo de Zhang--Shasha tem ainda como principal influência o trabalho de \citeonline{cooking_2025}, em que a recombinação de árvores por distância de edição gera ideias mais diversas e novas do que a geração direta por modelos de linguagem. Essa evidência, obtida em outro domínio, motiva investigar se a mesma família de métodos favorece a diversidade na geração de moléculas, mas não a transfere diretamente: o uso aqui é diferente (Seção~\ref{sec:estado-da-arte}), e a comparação com o Tanimoto e com a seleção aleatória, nas hipóteses H1 e H2, é o que permitirá dizer se a TED traz algum ganho.

\section{Limitações reconhecidas}

A TED mede a distância entre as representações estruturais, e não entre espectros, e a distância entre dois anéis pode não refletir a distância entre as moléculas finais que os contêm (hipótese H2). O indicador de conjugação usado como guia é uma triagem, e não uma previsão do comprimento de onda de absorção. Ambos os pontos são tratados como hipóteses a testar, e não como premissas.
